%% ======================================================================
%% appD/appendix_d_main.tex
%%
%% Appendix D: Security and Coverage Findings
%%
%% Expanded narratives and reference tables supporting Chapter 7's
%% security and user-centric results. Covers the two remediated security
%% vulnerabilities (Finding B authorization gap, Finding E view-pack
%% overwrite), the three accepted design properties (Findings A, C, D),
%% the full error-variant test coverage matrix, and the breakdown of
%% cargo audit advisories from the transitive Polkadot SDK dependencies.
%%
%% Sections:
%%   D.1 Security Audit Finding B: Authorization Gap
%%       in xcm_transfer_ownership
%%   D.2 Security Audit Finding E: View Pack Overwrite
%%   D.3 Accepted Design Findings (A, C, D)
%%   D.4 Error Variant Coverage (Table: 19 error variants, 16 tested)
%%   D.5 Static Analysis Advisory Breakdown (Table: 8 advisories)
%%
%% External labels referenced: none (self-contained appendix).
%%
%% Citations: none.
%% ======================================================================

\chapter{Security and Coverage Findings}\label{app:security}

\section[Finding B: Authorization Gap]{Security Audit Finding B:\texorpdfstring{\\}{ }Authorization Gap in \texttt{xcm\_transfer\_ownership}}\label{sec:d-finding-b}

The original implementation performed only \texttt{ensure\_signed(origin)} and then looked up \texttt{Ownership[content\_id, from]}, burned the seller's NFT, and minted a new one for \texttt{to}. No check confirmed that the dispatch caller had any relationship to the \texttt{from} account, so any signed account could call the extrinsic with arbitrary \texttt{from} and \texttt{to} arguments and transfer ownership from any user to itself.

We rate this finding \textbf{High}. Exploitation requires only a signed account with sufficient balance for transaction fees, no relationship to the victim, and no prior state; the impact is unauthorized transfer of a permanent, transferable ownership right. We initially recorded the finding as Medium on the grounds that the affected extrinsic is reachable only through the cross-chain path in normal operation, but Finding~A establishes that the \texttt{xcm\_*} extrinsics are also callable locally by any signed account, which removes that mitigation.

Our remediation was \texttt{ensure!(authorizer == from,\allowbreak\ Error::<T>::Unauthorized)}, committed alongside unit test \texttt{security\_xcm\_transfer\_\allowbreak ownership\_\allowbreak any\_account\_\allowbreak can\_steal}.

This remediation is deliberately conservative. Because an XCM \texttt{Transact} with \texttt{SovereignAccount} origin dispatches under the sending parachain's sovereign account, the equality check cannot be satisfied by a remote caller, so \texttt{xcm\_transfer\_ownership} no longer supports genuine cross-chain transfers and is equivalent to \texttt{transfer\_ownership} invoked locally. We accepted the loss of functionality in preference to a delegated-authority model that we could not verify within the implementation window; Appendix~\ref{sec:e-delegation} describes the capability-based alternative.

\section{Security Audit Finding E: View Pack Overwrite}\label{sec:d-finding-e}

The original \texttt{purchase\_views} always overwrote \texttt{ViewPacks[content\_id, buyer]} with a new \texttt{ViewPackInfo}, so a user holding 10 unused views who purchased 5 more would end up with 5, not 15. Our remediation made purchases additive: if a view pack already exists, new views are added to the existing balance.

This remediation covers the local \texttt{purchase\_views} path only. The XCM path (\texttt{xcm\_purchase\_views}) still overwrites \texttt{ViewPacks} unconditionally and mints a fresh child NFT on each call, orphaning the previous child. Extending the fix to the XCM path is deferred to future work.

\section{Accepted Design Findings (A, C, D)}\label{sec:d-accepted}

\textbf{Finding A:} \texttt{xcm\_*} extrinsics are callable locally as gift transactions (economically neutral; documented for transparency). \textbf{Finding C:} Runtime XCM \texttt{SafeCallFilter} is set to \texttt{Everything}; the \texttt{XcmExecuteFilter}, \texttt{XcmTeleportFilter}, and \texttt{XcmReserveTransferFilter} are similarly set to \texttt{Everything}, and \texttt{FixedWeightBounds} is used in place of measured XCM weights; production should restrict all four to explicit allowlists. \textbf{Finding D:} Content can be registered with all prices at zero (legitimate free-content feature, not a vulnerability).

\section{Finding I: Zero-View Denial of Service}\label{subsec:finding-i}

\textbf{Severity:} Medium

\textbf{Description.} \texttt{purchase\_views} accepts \texttt{num\_views = 0}, charges \texttt{ppv\_price} $\times$ 0 = 0, and still mints a child NFT. \texttt{consume\_view} then fails with \texttt{NoViewsRemaining}, so the slot is never released. Fifty throwaway accounts can permanently block all future sales of any content item for the cost of transaction fees alone.

\textbf{Impact.} Denial of service: any user can exhaust the fifty-child cap on any content item, preventing all future purchases at negligible cost. This is the exploit path for the fifty-child bound discussed in Section~\ref{sec:limitations}.

\textbf{Remediation.} Add \texttt{ensure!(num\_views > 0, Error::<T>::InsufficientPayment)} at the start of \texttt{purchase\_views} and \texttt{xcm\_purchase\_views}. A dedicated \texttt{ZeroViews} error variant would improve clarity.

\textbf{Status:} Open.

\section{Error Variant Coverage}\label{sec:d-coverage}

\begin{longtable}{@{}L{0.34\linewidth} L{0.12\linewidth} L{0.34\linewidth} L{0.16\linewidth}@{}}
\caption{Error Variant Coverage}\label{tab:error-coverage}\\
\toprule
\textbf{Error Variant} & \textbf{Tested} & \textbf{Error Variant} & \textbf{Tested} \\
\midrule
\endfirsthead

\multicolumn{4}{l}{\emph{Table~\ref{tab:error-coverage} continued from previous page}}\\
\toprule
\textbf{Error Variant} & \textbf{Tested} & \textbf{Error Variant} & \textbf{Tested} \\
\midrule
\endhead

\midrule
\multicolumn{4}{r}{\emph{Continued on next page}}\\
\endfoot

\bottomrule
\endlastfoot

\texttt{ContentNotFound} & Yes & \texttt{MaxChildrenReached} & Yes \\
\texttt{NotContentCreator} & Yes & \texttt{ContentIdOverflow} & Yes \\
\texttt{SubscriptionAlreadyExists} & Yes & \texttt{ItemIdOverflow} & No \\
\texttt{SubscriptionNotFound} & Yes & \texttt{OwnershipNotFound} & Yes \\
\texttt{SubscriptionNotExpired} & Yes & \texttt{InvalidRoyaltySplits} & Yes \\
\texttt{ViewPackNotFound} & Yes & \texttt{AutoRenewAlreadyEnabled} & Yes \\
\texttt{NoViewsRemaining} & Yes & \texttt{AutoRenewNotEnabled} & Yes \\
\texttt{AlreadyOwned} & Yes & \texttt{Unauthorized} & Yes \\
\texttt{InsufficientPayment} & Yes & \texttt{NftOperationFailed} & No \\
                                &     & \texttt{ChildAlreadyNested} & No (unreachable) \\
\end{longtable}

\section{Static Analysis Advisory Breakdown}\label{sec:d-static-analysis}

\begin{longtable}{@{}L{0.18\linewidth} L{0.24\linewidth} L{0.22\linewidth} L{0.30\linewidth}@{}}
\caption{Static Analysis Advisory Breakdown}\label{tab:advisories}\\
\toprule
\textbf{Category} & \textbf{Package} & \textbf{Risk} & \textbf{Mitigation} \\
\midrule
\endfirsthead

\multicolumn{4}{l}{\emph{Table~\ref{tab:advisories} continued from previous page}}\\
\toprule
\textbf{Category} & \textbf{Package} & \textbf{Risk} & \textbf{Mitigation} \\
\midrule
\endhead

\midrule
\multicolumn{4}{r}{\emph{Continued on next page}}\\
\endfoot

\bottomrule
\endlastfoot

Networking (2) & \texttt{quinn-proto}, \texttt{ring} & DoS, panic & Relay/p2p layer only; no runtime impact \\
\addlinespace
WASM executor (4) & \texttt{wasmtime} & Resource exhaustion, panics, f64 segfault & Sandboxed WASM execution model \\
\addlinespace
TLS (1) & \texttt{rustls-webpki} & CRL bypass & Node-side only \\
\addlinespace
Logging (1) & \texttt{tracing-subscriber} & Log injection & Observer layer only \\
\end{longtable}

All 8 advisories are mitigated by the architectural separation between the WASM-sandboxed runtime and the node binary. None are in code we authored.
